% TODO make hw stylesheet
\documentclass[10pt]{article}
\usepackage[letterpaper,top=1in,left=1in,right=1in]{geometry}
\usepackage[dvipsnames,svgnames]{xcolor}
\usepackage[shortlabels]{enumitem}
\usepackage[framemethod=TikZ]{mdframed}
\usepackage{amsmath,amssymb,amsthm}
\usepackage[colorlinks]{hyperref}
\usepackage{multicol}
\usepackage{tabularx}
\usepackage{bbm}

\hypersetup{
    colorlinks=true,
    linkcolor=blue,
    filecolor=magenta,
    urlcolor=magenta,
    pdftitle={MAT 344 HW}
}

\usepackage{mathtools}
\usepackage{thmtools}
\usepackage[textsize=footnotesize]{todonotes}
\usepackage{titling}
\usepackage{cancel}

\reversemarginpar
\mdfdefinestyle{mdbluebox}{roundcorner=10pt,innerbottommargin=9pt,
linecolor=blue,backgroundcolor=TealBlue!5,}
\declaretheoremstyle[headfont=\sffamily\bfseries\color{MidnightBlue},
mdframed={style=mdbluebox},]{thmbluebox}

\mdfdefinestyle{mdbloobox}{frametitlefont=\bfseries,innerbottommargin=8pt,
nobreak=true,backgroundcolor=TealBlue!5,linecolor=blue,}
\declaretheoremstyle[headfont=\bfseries\color{MidnightBlue},
mdframed={style=mdbloobox},headpunct={\\[3pt]},postheadspace=0pt,]{thmbloobox}

\mdfdefinestyle{mdredbox}{frametitlefont=\bfseries,innerbottommargin=8pt,
nobreak=true,backgroundcolor=Salmon!5,linecolor=RawSienna,}
\declaretheoremstyle[headfont=\bfseries\color{RawSienna},
mdframed={style=mdredbox},headpunct={\\[3pt]},postheadspace=0pt,]{thmredbox}

\mdfdefinestyle{mdgreenbox}{linecolor=ForestGreen,backgroundcolor=ForestGreen!5,
linewidth=2pt,rightline=false,leftline=true,topline=false,bottomline=false,}
\declaretheoremstyle[headfont=\bfseries\sffamily\color{ForestGreen!70!black},
mdframed={style=mdgreenbox},headpunct={ --- },]{thmgreenbox}

\mdfdefinestyle{mdorangebox}{linecolor=RedOrange,backgroundcolor=RedOrange!5,
linewidth=2pt,rightline=false,leftline=true,topline=false,bottomline=false,}
\declaretheoremstyle[headfont=\bfseries\sffamily\color{RedOrange!70!black},
mdframed={style=mdorangebox},headpunct={ --- },]{thmorangebox}

\mdfdefinestyle{mdblackbox}{linecolor=black,backgroundcolor=RedViolet!5!gray!5,
linewidth=3pt,rightline=false,leftline=true,topline=false,bottomline=false,}
\declaretheoremstyle[mdframed={style=mdblackbox}]{thmblackbox}

\declaretheoremstyle[spaceabove=6pt,spacebelow=6pt]{boldhead}
\declaretheoremstyle[spaceabove=6pt,spacebelow=6pt,headfont=\normalfont\itshape]{italichead}

\declaretheorem[style=boldhead,name=Problem]{problem}
\declaretheorem[style=boldhead,name=Problem,numbered=no]{problem*}
\declaretheorem[style=boldhead,name=Setup]{setup} % for config geo
\declaretheorem[style=boldhead,name=Example,sibling=problem]{example}
\declaretheorem[style=boldhead,name=Theorem,sibling=problem]{theorem}
\declaretheorem[style=boldhead,name=Corollary,sibling=theorem]{corollary}
\declaretheorem[style=boldhead,name=Proposition,sibling=theorem]{prop}
\declaretheorem[style=boldhead,name=Lemma,numberwithin=problem]{lemma}
\declaretheorem[style=boldhead,name=Theorem,numbered=no]{theorem*}
\declaretheorem[style=boldhead,name=Lemma,numbered=no]{lemma*}
\declaretheorem[style=boldhead,name=Claim,numberwithin=problem]{claim}
\declaretheorem[style=boldhead,name=Claim,numbered=no]{claim*}
\declaretheorem[style=boldhead,name=Remark,numberwithin=problem]{remark}
\declaretheorem[style=boldhead,name=Remark,numbered=no]{remark*}
\declaretheorem[style=boldhead,name=Definition,sibling=theorem]{definition}
\declaretheorem[style=boldhead,name=Definition,numbered=no]{definition*}

\providecommand{\ol}{\overline}
\providecommand{\ceil}[1]{\left\lceil #1 \right\rceil}
\providecommand{\floor}[1]{\left\lfloor #1 \right\rfloor}
\newcommand{\dd}{\mathrm{d}}
\providecommand{\ul}{\underline}
\providecommand{\eps}{\varepsilon}
\providecommand{\half}{\frac{1}{2}}
\providecommand{\CC}{\mathbb C}
\providecommand{\FF}{\mathbb F}
\providecommand{\II}{\mathbb I}
\providecommand{\NN}{\mathbb N}
\providecommand{\QQ}{\mathbb Q}
\providecommand{\RR}{\mathbb R}
\providecommand{\ZZ}{\mathbb Z}
\providecommand{\ind}{\mathbbm 1}
\providecommand{\dg}{^\circ}
\providecommand{\ii}{\item}
\providecommand{\setdiff}{\smallsetminus}
\providecommand{\alert}[1]{{\sffamily\textbf{\textcolor{blue}{#1}}}}
\providecommand{\inv}{^{-1}}
\providecommand{\mb}[1]{\mathbf{#1}}
\newcommand{\what}{\widehat}

\newenvironment{soln}{\begin{proof}[Solution]}{\end{proof}}

\providecommand{\pow}{\operatorname{Pow}}
\providecommand{\vocab}[1]{{\textbf{\textcolor{ForestGreen}{#1}}}}
\providecommand{\lcm}{\operatorname{lcm}}
\providecommand{\abs}[1]{\left\lvert #1\right\rvert}
\providecommand{\smabs}[1]{\lvert #1\rvert}
\providecommand{\innerprod}[1]{\langle\!\langle #1\rangle\!\rangle}
\providecommand{\norm}[1]{\left\| #1\right\|}
\providecommand{\smnorm}[1]{\| #1\|}
\providecommand{\gr}{\operatorname{gr}}
\providecommand{\im}{\operatorname{im}}

\usepackage{mathrsfs}

% place title at top right
\setlength{\droptitle}{-8ex}
\pretitle{\begin{flushright}\Large\bfseries}
\posttitle{\par\end{flushright}}
\preauthor{\begin{flushright}}
\postauthor{\end{flushright}}
\predate{\begin{flushright}}
\postdate{\end{flushright}}

\title{MAT 344 HW05}
\author{\large Aditya Pahuja}
\date{\large Due 03/03}
\begin{document}
\maketitle

\begin{problem}
    How many roots does the equation
    $z^7 - 2z^5 + 6z^3 - z + 1 = 0$ have
    in the disk $|z| < 1$?
\end{problem}

\begin{soln}
    For $\abs z = 1$,
    \begin{equation*}
	\abs{z^7 - 2z^5 - z + 1} \le 5 < 6 = \abs{6z^3},
    \end{equation*}
    so Rouch\'e's theorem implies that
    $z^7 - 2z^5 + 6z^3 - z + 1 = 0$ has the same number of solutions
    in the unit disk as $6z^3 = 0$, namely $\boxed 3$.
\end{soln}

\begin{problem}
    How many roots of the equation $z^4 - 6z + 3 = 0$
    have their modulus between $1$ and $2$?
\end{problem}

\begin{soln}
    Observe that
    \begin{equation*}
	\abs{-6z + 3} \le 6\cdot 2 + 3 = 15 < 16 = \abs{z^4}
    \end{equation*}
    when $\abs z = 2$, so the given equation has all four roots
    in the disk centered at $0$ with radius $2$.
    On the other hand,
    \begin{equation*}
	\abs{z^4 + 3} \le 4 < 6 = \abs{6z}
    \end{equation*}
    when $\abs z = 1$, so the given equation has one root inside
    (or on the boundary of) the disk centered at $0$ with radius $1$.
    Therefore $4-1=\boxed{3}$ of the solutions to the given equation
    have modulus between $1$ and $2$.
\end{soln}

\begin{problem}
    How many roots of the equation
    $z^4 + 8z^3 + 3z^2 + 8z + 3 = 0$
    lie in the right half-plane?
\end{problem}

\begin{soln}
    Let $R > 0$ be a real number.
    The image under the given polynomial $f$
    of the line segment joining $iR$ and $-iR$
    is contained in the right half-plane because for any $r\in\RR$,
    $f(ir) = (r^4 - 3r^2 + 3) + i(-8r^3 + 8r)$
    and $r^4 - 3r^2 + 3 = (r^2 - \frac32)^2 + \frac34 > 0$
    when $r$ is real, so this image is path homotopic to
    the segment joining $f(iR)$ and $f(-iR)$.
    This means that the winding number of the image under $f$ of
    the loop formed by the segment joining $-iR$ and $iR$
    and the counterclockwise semicircular arc from $-iR$ to $iR$
    is the same as the winding number of the loop formed by
    the image under $f$ of the semicircular arc and the segment
    joining $f(iR)$ to $f(-iR)$.

    Parametrize the semicircular arc as a function on $[-\pi/2,\pi/2]$
    in the obvious way and parametrize the segment from $f(iR)$ to $f(-iR)$
    as a function on $[0,\pi]$; the concatenation of these two paths
    is a loop which spends half its time traversing the semicircle
    and half its time traversing the segment.
    Reparametrize the loop by scaling and shifting its domain to be $[0,1]$,
    and call the loop $\gamma\colon[0,1]\to\CC-\{0\}$
    (to be precise, $\gamma$ traverses the semicircle on $[0,1/2]$
    and the segment on $[1/2,1]$); recall that our goal is
    to compute the winding number of $\gamma$ around $0$.
    Let $\omega\colon[0,1]\to\CC-\{0\}$ be the loop
    $\omega(t) = R^4e^{8\pi it}$ for $0\le t\le 1/2$
    and $\omega(t) = R^4$ for $1/2\le t\le 1$, so that
    $\omega$ winds around $0$ twice before staying for a while at $R^4$.
    Then,
    \begin{equation*}
	\abs{\gamma(t) - \omega(t)} < \abs{\omega(t)} = R^4
    \end{equation*}
    for large enough $R$, noting that
    $f(\pm iR) - R^4 = (-3R^2 + 3) \pm i(-8R^3 + 8R)$
    is only cubic in $R$ so that the inequality holds for $t\in[1/2,1]$.
    It follows from a similar argument to Rouch\'e's theorem that
    $\gamma$ and $\omega$ have the same winding number of \boxed{2}.

    To be thorough, the inequality tells us that
    for each $t$, the segment joining $\gamma(t)$ to $\omega(t)$
    does not contain $0$, so there is a homotopy
    $H_s(t)\colon[0,1]\times[0,1]\to\CC-\{0\}$
    such that for a given $t$, $H_s(t)$ is the segment from $\gamma(t)$
    to $\omega(t)$. This map has the property that for each $s$,
    $H_s(0) = H_s(1)$, so $H$ descends to a homotopy
    $S^1\times[0,1]\to\CC-\{0\}$ via the quotient map
    $[0,1]\to S^1$. Such a homotopy preserves winding numbers
    because the change of basepoint isomorphism
    from $\pi_1(\CC-\{0\},\gamma(0))$ to $\pi_1(\CC-\{0\},\omega(0))$
    along the path $H_s(0)\colon[0,1]\to\CC-\{0\}$
    then sends the homotopy class of $\gamma$
    to the homotopy class of $\omega$.
\end{soln}

% \begin{proof}
%     [Solution 2]
%     In fact we can explicitly solve the given polynomial:
%     $(z^2 + 1)^2 - 2z^2 + 8z(z^2 + 1) + 3z^2
% \end{proof}










\end{document}
